% PCT-SOURCE: pdf=180 print=168 kind=section-title id=sec-ch4-6-borchers-classes
\subsection{Equivalence Classes of Local Fields (Borchers Classes)}
\label{sec:ch4-6-borchers-classes}

% PCT-SOURCE: pdf=180 print=168 kind=prose id=p168-borchers-introduction
One can arrive at the notion of equivalence classes of local fields by
several routes. The most direct is the following. Let $\varphi_1$ be a
hermitian scalar field with cyclic vacuum state. Suppose $\varphi_2$ and
$\varphi_3$ are two other fields with the same representation of the
Poincaré group as $\varphi_1$. Suppose $\varphi_2$ and $\varphi_3$ are not
necessarily local but instead are \emph{local relative} to $\varphi_1$. That
means

% PCT-SOURCE: pdf=180 print=168 kind=display id=4-90
\begin{equation}
\begin{aligned}
  [\varphi_1(x),\varphi_2(y)]_-&=0,\\
  [\varphi_1(x),\varphi_3(y)]_-&=0,
  \qquad \text{for }(x-y)^2>0.
\end{aligned}
\tag{4-90}\label{eq:ch4-relative-locality-assumptions}
\end{equation}

% PCT-SOURCE: pdf=180 print=168 kind=prose id=p168-borchers-discovery
Then what can be said about the relation of $\varphi_2$ to $\varphi_3$?
Borchers discovered that they are relatively local:

% PCT-SOURCE: pdf=180 print=168 kind=display id=4-91
\begin{equation}
  [\varphi_2(x),\varphi_3(y)]_-=0,
  \qquad \text{for }(x-y)^2>0.
\tag{4-91}\label{eq:ch4-relative-locality-conclusion}
\end{equation}

% PCT-SOURCE: pdf=180 print=168 kind=prose id=p168-borchers-equivalence
That immediately implies that they are local (take $\varphi_2=\varphi_3$).
Thus, relatively local fields fall into equivalence classes.\footnote{%
% PCT-SOURCE: pdf=180 print=168 kind=footnote id=fn-p168-borchers-class
We should mention at the outset that what is proved is that if a local
$\varphi_1$ has the vacuum as cyclic vector, which implies that $\varphi_1$
is irreducible, then (4-90) implies (4-91). The $\varphi_2$, $\varphi_3$ need
not be irreducible. Now in the customary mathematical definition of
equivalence all elements of an equivalence class must be on the same footing,
so strictly speaking it is the \emph{irreducible} relatively local fields
which form equivalence classes. In the following, we shall use the term
equivalence class somewhat loosely, including reducible relatively local
fields in the equivalence class along with the irreducible.}

% PCT-SOURCE: pdf=181 print=169 kind=prose id=p169-borchers-purpose
The main purpose of this section is to prove these results and analogous ones
involving weak local commutativity. However, before we begin, we want to give
some idea of their significance.

% PCT-SOURCE: pdf=181 print=169 kind=prose id=p169-wick-polynomials
What do the fields look like which lie in a given equivalence class? In
Section 3-2 we have introduced the notion of a Wick polynomial in a free
field. It can be shown that the Wick polynomials in a given free field exhaust
the equivalence class of that free field (Ref.~27). It is plausible for fields
which are not free that the relations between fields in the same equivalence
class will be something like those between the Wick polynomials; they can be
expressed as local functions of each other. An example is $\varphi(x)$ and
$\varphi(x)+(-\Box+m^2)\varphi(x)$. At the moment, the investigation of the
precise meaning which should be given to the idea that one field is a local
function of another is under intensive investigation but the matter is
scarcely in definitive shape. It is closely connected with the theory of von
Neumann algebras of open sets mentioned at the end of Section 4-2.

% PCT-SOURCE: pdf=181 print=169 kind=prose id=p169-s-matrix-significance
The physical importance of the equivalence relation that arises from relative
local commutativity is this: two field theories which have the same Hilbert
space $\Hilbert$ and transformation law $U$ and whose fields are relatively
local have the same $S$-matrix. The proof of this statement makes use of the
Haag--Ruelle theory and so is outside the scope of this book. Nevertheless,
because of its importance we want to give some idea of the direction in which
the argument goes and at the same time show how equivalence can be practically
important.

% PCT-SOURCE: pdf=181 print=169 kind=prose id=p169-haag-ruelle-example
Consider the theory of proton, neutron, and $\pi$-mesons, $\pi^\pm$ and
$\pi^0$, already discussed at the end of Section 4-2, in which the proton and
neutron fields, $\psi_p$, $\psi_n$, together with their adjoints form an
irreducible set of operators. In the construction of Haag and Ruelle, one
forms states $Q_\alpha(t)\ket{\Omega}$, where $Q_\alpha(t)$ is a suitably
smeared polynomial in these fields such that

% PCT-SOURCE: pdf=181 print=169 kind=display id=p169-collision-limit
\begin{equation*}
  \lim_{t\to+\infty}
  \left\lVert Q_\alpha(t)\ket{\Omega}-\OutKet{\Psi_\alpha}\right\rVert=0.
\end{equation*}

% PCT-SOURCE: pdf=181 print=169 kind=prose id=p169-collision-states
where the $\InKet{\Psi_\alpha}$ are the collision states of neutrons, protons,
and $\pi$ mesons. In this theory, there exists another irreducible set of
operators, $\psi_n$, $\varphi^+$, and their adjoints which can create all the
states. These fields can also be used to obtain collision states using the
Haag--Ruelle prescription. Ruelle's theory assures us that if the fields
$\psi_p$, $\psi_n$, $\varphi^+$, $\varphi^-$, $\varphi^0$ are local and local
relative to one another, then the collision states defined in terms of one
irreducible set, $\psi_p$, $\psi_n$, are the same as the collision states
defined in terms of another $\psi_n$, $\varphi^+$. It follows that the
$S$-matrix is the same whichever set is used for the definition of the
asymptotic states.

% PCT-SOURCE: pdf=182 print=170 kind=prose id=p170-s-equivalence-problem
One gets further perspective on the significance of the equivalence classes
of local fields if one considers the following problem: classify field
theories with a given $\Hilbert$ and $U$ which have different $S$-matrices.
(A really effective solution of this problem would, of course, be a crowning
achievement of the general theory of quantized fields even if it did not lead
to the calculation of a single cross section.) Suppose we call two field
theories $S$-equivalent if they lead to the same $S$-matrix. Then the
preceding can be restated: two field theories in $\Hilbert$, with the same $U$
and relatively local fields, are $S$-equivalent. Thus, $S$-equivalence classes
are made up of Borchers classes. It is not difficult to see that the same
$S$-equivalence class will contain many Borchers classes. For example, from a
field theory of a neutral scalar field $\varphi$, one can obtain new theory by
replacing $\varphi$ by $V\varphi V^{-1}$, where $V$ is any unitary
transformation commuting with the representation $U$ of
$\mathcal P_+^\uparrow$. In general, $V\varphi V^{-1}$ will not be local
relative to $\varphi$, but it will yield an $S$-equivalent theory. It is not
yet known whether all the field theories of a given $S$-equivalence class can
be obtained from those of a single Borchers' class by applying all possible
$V$'s in this way. What is clear is that a knowledge of the structure of
Borchers classes is an essential preliminary in any effort to understand the
structure of $S$-equivalence classes.

% PCT-SOURCE: pdf=182 print=170 kind=prose id=p170-wlc-definition
With these qualitative remarks completed, we now commence our treatment of
equivalence classes. For simplicity, we will discuss field theories given in
terms of a single hermitian scalar field. The generalization to arbitrary sets
of spinor fields is straightforward. We recall that a field, $\varphi$, is
weakly local if the equality

% PCT-SOURCE: pdf=182 print=170 kind=display id=4-92
\begin{equation}
  \bra{\Omega}\varphi(x_1)\cdots\varphi(x_n)\ket{\Omega}
  =
  \bra{\Omega}\varphi(x_n)\cdots\varphi(x_1)\ket{\Omega}
  \tag{4-92}\label{eq:ch4-wlc}
\end{equation}

% PCT-SOURCE: pdf=182 print=170 kind=prose id=p170-wlc-jost
holds for each $n$ and each Jost point $x_1\ldots x_n$. We saw in the
\emph{PCT} theorem that $\varphi$ is weakly local if and only if there exists
an anti-unitary $\Theta$ satisfying

% PCT-SOURCE: pdf=182 print=170 kind=display id=4-93
\begin{equation}
  \Theta\varphi(f)\Theta^{-1}=\varphi(\widehat f).
  \tag{4-93}\label{eq:ch4-wlc-pct}
\end{equation}

% PCT-SOURCE: pdf=182 print=170 kind=display id=4-93-conjugate-test-function
\begin{equation*}
  \widehat f(x)=f^*(-x).
\end{equation*}

% PCT-SOURCE: pdf=182 print=170 kind=prose id=p170-pct-uniqueness
Since by Theorem 4-5 $\varphi$ is irreducible, (4-93) determines $\Theta$ up
to a phase (the proof is in\textsuperscript{\(\dagger\)} Section 3-5). If there
is another irreducible field $\psi$ defined in $\Hilbert$ with the same
representation $U$ of $\mathcal P_+^\uparrow$, it will have a \emph{PCT}
operator $\Theta_\psi$, say, if and only if it is weakly local. We then have
two \emph{PCT} operators $\Theta$ and $\Theta_\psi$; when will they be the
same? The reconstruction theorem (Theorem 3-7) and Theorem 4-7 show that if
the domains of the two fields satisfy certain conditions, and if $\varphi$ and
$\psi$ together satisfy the \emph{WLC} identities, then
$\Theta=\Theta_\psi$. We shall describe

% PCT-SOURCE: pdf=183 print=171 kind=prose continuation-of=p170-pct-uniqueness
this situation by saying the fields are \emph{weakly local with respect to
each other}. As an example of a field theory with three irreducible fields
which are not weakly local with respect to each other, consider a non-trivial
theory of a hermitian scalar field $\varphi(x)$ which has asymptotic free
fields $\varphi^{\mathrm{in}}(x)$, $\varphi(x)$, $\varphi^{\mathrm{out}}(x)$,
and suppose the theory is asymptotically complete. Since

% PCT-SOURCE: pdf=183 print=171 kind=display id=p171-cpt-in-out
\begin{equation*}
  \Theta\varphi^{\mathrm{in}}(f)\Theta^{-1}
  =\varphi^{\mathrm{out}}(\widehat f)
  \ne \varphi^{\mathrm{in}}(\widehat f).
\end{equation*}

% PCT-SOURCE: pdf=183 print=171 kind=prose id=p171-cpt-distinction
the \emph{CPT} operator $\Theta$ for $\varphi(x)$ is not the \emph{CPT}
operator for $\varphi^{\mathrm{in}}$ or $\varphi^{\mathrm{out}}$. It follows
\emph{a fortiori} that $\varphi^{\mathrm{in}}(x)$, $\varphi(x)$, and
$\varphi^{\mathrm{out}}(x)$ cannot be local with respect to each other unless
all the fields coincide.

% PCT-SOURCE: pdf=183 print=171 kind=prose id=p171-wlc-conditions
We now give a set of conditions that $\varphi$ and $\psi$ be weakly local
relative to each other. These conditions are a subset of those which occur in
Theorem 4-7. It should be noted that neither field $\varphi$ nor $\psi$ of the
theorem is assumed local, i.e., neither need satisfy axiom III.

% PCT-SOURCE: pdf=183 print=171 kind=theorem id=4-18
\begin{theorem}[4-18]\label{thm:ch4-18-weak-local-relative}
% PCT-SOURCE: pdf=183 print=171 kind=theorem-statement id=4-18-statement
Suppose $\varphi(x)$ is a weakly local field with vacuum as cyclic vector, and
suppose $\psi(x)$ is another field transforming under the same representation
of $\mathcal P_+^\uparrow$, and with the same domain $D$ of definition.
Suppose

% PCT-SOURCE: pdf=183 print=171 kind=display id=4-94
\begin{equation}
\begin{aligned}
 &\bra{\Omega}\varphi(x_1)\cdots\varphi(x_j)\psi(x)
   \varphi(x_{j+1})\cdots\varphi(x_n)\ket{\Omega}\\
 &\quad =
   \bra{\Omega}\varphi(x_n)\cdots\varphi(x_{j+1})\psi(x)
   \varphi(x_j)\cdots\varphi(x_1)\ket{\Omega}
\end{aligned}
\tag{4-94}\label{eq:ch4-wlc-interchange}
\end{equation}

% PCT-SOURCE: pdf=183 print=171 kind=theorem-statement-continuation id=4-18-conclusion
holds at Jost points for all $j$ and $n$; then
\begin{enumerate}
\item[(a)] $\psi(x)$ is a weakly local field, and
\item[(b)] $\varphi(x)$, $\psi(x)$ are weakly local with respect to each other.
\end{enumerate}
\end{theorem}

% PCT-SOURCE: pdf=183 print=171 kind=proof id=proof-ch4-18
\begin{proof}
% PCT-SOURCE: pdf=183 print=171 kind=prose id=p171-proof-4-18-pct
If $\Theta$ is the \emph{PCT} operator for the field $\varphi$, we have for
any states $\ket{\Phi}$, $\ket{\Psi}\in D$

% PCT-SOURCE: pdf=183 print=171 kind=display id=4-95
\begin{equation}
  \matrixel{\Theta\Phi}{\Theta\psi(x)\Theta^{-1}}{\Theta\Psi}
  =
  \matrixel{\Phi}{\psi(x)}{\Psi}^{*}.
  \tag{4-95}\label{eq:ch4-antiunitary-field}
\end{equation}

% PCT-SOURCE: pdf=183 print=171 kind=prose id=p171-proof-4-18-continuation
Using the invariance of (4-94) under the complex Lorentz group we get for
points of holomorphy

% PCT-SOURCE: pdf=183 print=171 kind=display id=4-96
\begin{equation}
\begin{aligned}
 &\bra{\Omega}\varphi(x_1)\cdots\varphi(x_j)\psi(x)
   \varphi(x_{j+1})\cdots\varphi(x_n)\ket{\Omega}\\
 &\quad -
   \bra{\Omega}\varphi(-x_n)\cdots\varphi(-x_{j+1})\psi(-x)
   \varphi(-x_j)\cdots\varphi(-x_1)\ket{\Omega}=0.
\end{aligned}
\tag{4-96}\label{eq:ch4-wlc-holomorphic}
\end{equation}

% PCT-SOURCE: pdf=183 print=171 kind=display id=4-97
\begin{equation}
  \operatorname{Im}(x_1-x_2),\ldots,\operatorname{Im}(x_j-x),
  \operatorname{Im}(x-x_{j+1}),\ldots,
  \operatorname{Im}(x_{n-1}-x_n)\in -V_+.
  \tag{4-97}\label{eq:ch4-wlc-tube}
\end{equation}

% PCT-SOURCE: pdf=183 print=171 kind=prose id=p171-proof-4-18-boundary
and so the equation holds at all points of holomorphy. Taking the limit from
the tube (4-97) we see that (4-96) holds at all real points. If

% PCT-SOURCE: pdf=183 print=171 kind=display id=4-98
\begin{equation}
  \ket{\Psi}
    =\varphi(f_1)\cdots\varphi(f_j)\ket{\Omega},
  \qquad
  \ket{\Phi}
    =\varphi(f_{j+1})\cdots\varphi(f_n)\ket{\Omega}.
  \tag{4-98}\label{eq:ch4-wlc-test-states}
\end{equation}

% PCT-SOURCE: pdf=184 print=172 kind=prose continuation-of=p171-proof-4-18-boundary
then

% PCT-SOURCE: pdf=184 print=172 kind=display id=p172-pct-test-state-images
\begin{equation*}
  \Theta\ket{\Psi}
    =\varphi(\widehat f_j)\cdots\varphi(\widehat f_1)\ket{\Omega},
  \qquad
  \Theta\ket{\Phi}
    =\varphi(\widehat f_{j+1})\cdots\varphi(\widehat f_n)\ket{\Omega}.
\end{equation*}

% PCT-SOURCE: pdf=184 print=172 kind=prose id=p172-proof-4-18-matrix-element
Using (4-96) we get

% PCT-SOURCE: pdf=184 print=172 kind=display id=p172-proof-4-18-comparison
\begin{equation*}
  \matrixel{\Psi}{\psi(x)}{\Phi}
  =\matrixel{\Theta\Phi}{\psi(-x)}{\Theta\Psi}.
\end{equation*}

% PCT-SOURCE: pdf=184 print=172 kind=prose id=p172-proof-4-18-conjugation
which when compared with (4-95) leads to

% PCT-SOURCE: pdf=184 print=172 kind=display id=p172-pct-psi
\begin{equation*}
  \Theta\psi(x)\Theta^{-1}=\psi(-x),
\end{equation*}

% PCT-SOURCE: pdf=184 print=172 kind=prose id=p172-proof-4-18-conclusion
an equation valid on states of the form (4-98). Using the hermiticity of
$\psi$, we conclude (just as in the argument at the end of Section 4-1) that
it is valid on all of $D$. Therefore there is a \emph{PCT} operator for the
field $\psi(x)$ [this proves (a)] and this operator is identical with the
\emph{PCT} operator of the field $\varphi(x)$ [this proves (b)].
\end{proof}

% PCT-SOURCE: pdf=184 print=172 kind=prose id=p172-transitivity-introduction
We now show that the property of being weakly local with respect to a given
irreducible field is \emph{transitive}, in a sense which the theorem states.
(The customary definition of transitivity for a relation $\equiv$ is
$a\equiv b$ and $b\equiv c$ imply $a\equiv c$. We have this strict
transitivity if $\varphi_2$ and $\varphi_3$ also have the vacuum as cyclic
vector.)

% PCT-SOURCE: pdf=184 print=172 kind=theorem id=4-19
\begin{theorem}[4-19]\label{thm:ch4-19-transitivity}
% PCT-SOURCE: pdf=184 print=172 kind=theorem-statement id=4-19-statement
Suppose $\varphi_1$ is a weakly local field for which the vacuum is cyclic,
and $\varphi_j$, $j=2,3$, are weakly local with respect to $\varphi_1$, and
have the same domain of definition and representation of $\mathcal
P_+^\uparrow$. Then $\varphi_2$ is weakly local with respect to $\varphi_3$.
\end{theorem}

% PCT-SOURCE: pdf=184 print=172 kind=proof id=proof-ch4-19
\begin{proof}
% PCT-SOURCE: pdf=184 print=172 kind=prose id=p172-proof-4-19
From the proof of Theorem 4-18 it follows that $\varphi_1$ and $\varphi_2$
have a simultaneous \emph{PCT} operator, and so have $\varphi_1$ and
$\varphi_3$. Since $\varphi_1$ is irreducible this operator is unique;
therefore $\varphi_2$ and $\varphi_3$ have the same \emph{PCT} operator. Thus
they are weakly local with respect to each other, by the \emph{PCT} theorem.
\end{proof}

% PCT-SOURCE: pdf=184 print=172 kind=prose id=p172-equivalence-class-result
The important result is that given a weakly local field $\varphi(x)$ for which
the vacuum is cyclic, we can form an equivalence class consisting of all
fields in the theory which are weakly local with respect to $\varphi$. All
fields with the same \emph{PCT} operator and domain are in this class. As far
as the $S$-equivalence of fields in this class is concerned we have the
following weak result.

% PCT-SOURCE: pdf=185 print=173 kind=theorem id=4-20
\begin{theorem}[4-20]\label{thm:ch4-20-asymptotic-fields}
% PCT-SOURCE: pdf=185 print=173 kind=theorem-statement id=4-20-statement
If $\varphi_1(x)$ is weakly local and has the vacuum as cyclic vector, and if
$\varphi_2(x)$ is weakly local with respect to $\varphi_1(x)$, and if there
are asymptotic fields such that

% PCT-SOURCE: pdf=185 print=173 kind=display id=4-99
\begin{equation}
  \varphi_1^{\mathrm{in}}(x)=\varphi_2^{\mathrm{in}}(x)
  \tag{4-99}\label{eq:ch4-asymptotic-in}
\end{equation}

% PCT-SOURCE: pdf=185 print=173 kind=theorem-statement-continuation id=4-20-then
then

% PCT-SOURCE: pdf=185 print=173 kind=display id=4-100
\begin{equation}
  \varphi_1^{\mathrm{out}}(x)=\varphi_2^{\mathrm{out}}(x).
  \tag{4-100}\label{eq:ch4-asymptotic-out}
\end{equation}

\end{theorem}

% PCT-SOURCE: pdf=185 print=173 kind=proof id=proof-ch4-20
\begin{proof}
% PCT-SOURCE: pdf=185 print=173 kind=prose id=p173-proof-4-20
$\varphi_1$ and $\varphi_2$ have the same \emph{PCT} operator, so since
$\Theta$ carries in-fields into out-fields, (4-100) follows from (4-99).
\end{proof}

% PCT-SOURCE: pdf=185 print=173 kind=prose id=p173-local-fields-stronger
For local fields there are stronger results, since, as remarked before,
(4-99) can be proved using the Haag--Ruelle theory. The next theorem asserts
that the relation of relative locality partitions local fields into classes.
This is the result whose significance was discussed at the beginning of this
section.

% PCT-SOURCE: pdf=185 print=173 kind=theorem id=4-21
\begin{theorem}[4-21]\label{thm:ch4-21-relative-locality-transitive}
% PCT-SOURCE: pdf=185 print=173 kind=theorem-statement id=4-21-statement
Suppose $\varphi_1$, $\varphi_2$, and $\varphi_3$ are fields with the same
domain and representation $U$ of $\mathcal P_+^\uparrow$, $\varphi_1$ is
local and has the vacuum as cyclic vector, and in addition

% PCT-SOURCE: pdf=185 print=173 kind=display id=p173-4-21-assumptions
\begin{equation*}
  [\varphi_1(x),\varphi_2(y)]=0=[\varphi_1(x),\varphi_3(y)]
  \qquad \text{if }(x-y)^2>0
\end{equation*}

% PCT-SOURCE: pdf=185 print=173 kind=theorem-statement-continuation id=4-21-conclusion
then

% PCT-SOURCE: pdf=185 print=173 kind=display id=p173-4-21-conclusion
\begin{equation*}
  [\varphi_2(x),\varphi_3(y)]=0
  \qquad \text{if }(x-y)^2>0.
\end{equation*}
\end{theorem}

% PCT-SOURCE: pdf=185 print=173 kind=proof id=proof-ch4-21
\begin{proof}
% PCT-SOURCE: pdf=185 print=173 kind=prose id=p173-proof-4-21-opening
The fields $\varphi_1$, $\varphi_2$, and $\varphi_3$ satisfy the conditions
of Theorem 4-18, and so $\varphi_2$ and $\varphi_3$ are weakly local, and
weakly local with respect to $\varphi_1$. Hence they are weakly local with
respect to each other. Thus at Jost points

% PCT-SOURCE: pdf=185 print=173 kind=display id=p173-4-21-jost
\begin{equation*}
\begin{aligned}
 &\bra{\Omega}\varphi_1(x_1)\cdots\varphi_1(x_j)
   \varphi_2(y)\varphi_3(z)\varphi_1(x_{j+1})\cdots
   \varphi_1(x_n)\ket{\Omega}\\
 &=
   \bra{\Omega}\varphi_1(x_n)\cdots\varphi_1(x_{j+1})
   \varphi_3(z)\varphi_2(y)\varphi_1(x_j)\cdots
   \varphi_1(x_1)\ket{\Omega}\\
 &=
   \bra{\Omega}\varphi_1(x_1)\cdots\varphi_1(x_j)
   \varphi_3(z)\varphi_2(y)\varphi_1(x_{j+1})\cdots
   \varphi_1(x_n)\ket{\Omega}.
\end{aligned}
\end{equation*}

% PCT-SOURCE: pdf=185 print=173 kind=prose id=p173-proof-4-21-commutation
since $\varphi_1$ is local, and commutes with $\varphi_2$ and $\varphi_3$.
This is precisely the relation studied in the proof of Theorem 4-1, and just
as there, we can conclude that for all $x_1\ldots x_n$,

% PCT-SOURCE: pdf=185 print=173 kind=display id=p173-4-21-vacuum-commutator
\begin{equation*}
  \bra{\Omega}\varphi_1(x_1)\cdots\varphi_1(x_j)
  [\varphi_2(y),\varphi_3(z)]
  \varphi_1(x_{j+1})\cdots\varphi_1(x_n)\ket{\Omega}=0
  \qquad \text{if }(y-z)^2>0.
\end{equation*}
\end{proof}

% PCT-SOURCE: pdf=186 print=174 kind=prose continuation-of=p173-proof-4-21
Since $\ket{\Omega}$ is cyclic for $\varphi_1$ we conclude that

% PCT-SOURCE: pdf=186 print=174 kind=display id=p174-4-21-final
\begin{equation*}
  [\varphi_2(y),\varphi_3(z)]=0
  \qquad \text{if }(y-z)^2>0,
\end{equation*}

% PCT-SOURCE: pdf=186 print=174 kind=prose id=p174-4-21-required
as required.

% PCT-SOURCE: pdf=186 print=174 kind=corollary id=corollary-ch4-21
\par\medskip\noindent\textbf{Corollary}\par\smallskip
% PCT-SOURCE: pdf=186 print=174 kind=corollary-statement id=corollary-ch4-21-statement
Putting $\varphi_2=\varphi_3$ we get, with the same assumptions:
If $\varphi_1(x)$ is local and has the vacuum as cyclic vector, and
$\varphi_2$ is local with respect to $\varphi_1$ then $\varphi_2$ is local.

% PCT-SOURCE: pdf=186 print=174 kind=prose id=p174-relative-locality-transitive
Theorem 4-21 shows that relative locality is a transitive property.

% PCT-SOURCE: pdf=186 print=174 kind=prose id=p174-uniqueness-application
As an application of Theorems 4-21 and 4-15 consider the problem of the
uniqueness of the solution of the equation

% PCT-SOURCE: pdf=186 print=174 kind=display id=p174-kg-equation
\begin{equation*}
  (-\Box+m^2)u(x)=j(x).
\end{equation*}

% PCT-SOURCE: pdf=186 print=174 kind=prose id=p174-uniqueness-setup
Here $j(x)$ is assumed to be a given irreducible local field, and the problem
is to find local $u$ which satisfy the equation. Suppose $u_1$ and $u_2$ were
two such. Since both are local relative to $j$ they are local relative to each
other. Thus their difference $\varphi=u_1-u_2$ is local and satisfies

% PCT-SOURCE: pdf=186 print=174 kind=display id=p174-free-equation
\begin{equation*}
  (\Box-m^2)\varphi(x)=0.
\end{equation*}

% PCT-SOURCE: pdf=186 print=174 kind=prose id=p174-free-field-argument
$\varphi$ is therefore a multiple of a free field. We should like to conclude
that $\varphi=0$. To make this step easy we add a further hypothesis. We
assume that the free field $u_{\mathrm{in}}$ exists,
$u_{\mathrm{in}}(f)\ket{\Omega}$ is a one-particle state $\ket{\Psi_{1f}}$,
and

% PCT-SOURCE: pdf=186 print=174 kind=display id=p174-one-particle-matrix
\begin{equation*}
  \bra{\Omega}u(x)\ket{\Psi_{1f}}
  =\bra{\Omega}u_{\mathrm{in}}(x)\ket{\Psi_{1f}},
\end{equation*}

% PCT-SOURCE: pdf=186 print=174 kind=prose id=p174-one-particle-test
for all $f\in\Schwartz$. It now follows that

% PCT-SOURCE: pdf=186 print=174 kind=display id=p174-difference-matrix
\begin{equation*}
  \bra{\Omega}[u_1(x)-u_2(x)]\ket{\Psi_{1f}}=0,
\end{equation*}

% PCT-SOURCE: pdf=186 print=174 kind=prose id=p174-uniqueness-conclusion
which in turn implies $\varphi=0$. Thus, we have

% PCT-SOURCE: pdf=186 print=174 kind=theorem id=4-22
\begin{theorem}[4-22]\label{thm:ch4-22-unique-local-solution}
% PCT-SOURCE: pdf=186 print=174 kind=theorem-statement id=4-22-statement
If $j$ is a given irreducible local field and $u(x)$ is a local solution of

% PCT-SOURCE: pdf=186 print=174 kind=display id=4-101
\begin{equation}
  (-\Box+m^2)u(x)=j(x)
  \tag{4-101}\label{eq:ch4-unique-local-solution}
\end{equation}

% PCT-SOURCE: pdf=186 print=174 kind=theorem-statement-continuation id=4-22-continuation
with a one-particle state at mass $m$, then $u$ is uniquely determined by the
requirement that

% PCT-SOURCE: pdf=186 print=174 kind=display id=p174-4-22-boundary-condition
\begin{equation*}
  \bra{\Omega}u(x)\ket{\Psi_{1f}}
  =\frac{1}{\ii}\int\Delta^+(x-y)\,\dd y\,f(y).
\end{equation*}
\end{theorem}

% PCT-SOURCE: pdf=186 print=174 kind=prose id=p174-4-22-final
There are other hypotheses under which the uniqueness of the local solution of
(4-101) follows but we shall not elaborate here.
